\section{Opening closed sums with local plugs}\label{sec:plugs}

The closed estimate \eqref{fields:closed} does not directly control a
proper subarc of a rim.  We shall close such subarcs in controlled local
experiments, compare the resulting closures, and eventually transfer the
comparison back to the original configuration.  The conclusion is the
following bulk statement.

\begin{theorem}[Bulk transport]\label{thm:bulk-transport}
From every sequence $\eps\downarrow0$ one can extract a subsequence and a
complex number $a$, depending only on that subsequence, such that
$\Rea a<1/2$ and the following holds in every bulk patch.  If
$f\in C_c^\infty(B(x,R))$ and $\overline B(x,6R)$ lies in the patch, then
\begin{equation}\label{transport:bulk}
 \max_{P:\,P\not\subset B(x,4R)}
       |Q_P(f)-aV_P(f)|\xrightarrow[\eps\to0]{\PP}0.
\end{equation}
As in the definition of a bulk patch, local tests may be set to zero on a
geometry event whose probability tends to zero.  The subsequence and
scalar are common to all patches and tests.
\end{theorem}

The transport estimate is completed in Section~\ref{sec:actual-arcs};
Proposition~\ref{prop:nondegeneracy} proves the strict inequality for $a$.
The output of this section is Proposition~\ref{prop:state-cochain}:
for each fixed connector, its contribution is asymptotically a function
of two retained states, and these functions have small cycle sums.
Here and in Section~\ref{sec:affine}, $X$ denotes either $Q$ or $V$.

\begin{convention}[The order of limits]\label{plug:scale-order}
A comparison scale $h>0$ is fixed before the microscopic limit.  Every
$o(1)$ and $o_{L^1}(1)$ in this section and the next is taken as
$\eps\downarrow0$ at fixed $h$ and fixed deterministic geometric pattern,
through values for which the stated types are admitted.  Later we use a
countable sequence $h\downarrow0$.  Thus a finite pattern may have an
arbitrarily small positive conditioning probability, and the number of
patterns may depend on $h$.  Bounds of the form $C h$ will instead be
obtained from fixed normalized short templates, with a common constant.
\end{convention}

\subsection{Types, states, roots and gates}

Normalized coordinates are multiplied by $h$ and placed in a chosen
oriented Euclidean frame.  We usually retain the same symbols for the
placed sets.  The following definition is for a black path; exchanging
the two colors gives the white types.

\begin{definition}[A plug type]\label{def:plug-type}
A type consists of the following deterministic data:
\begin{enumerate}[label=\textup{(\roman*)}]
\item two concentric gate annuli $[a_1,b_1]$, $[a_2,b_2]$ and a cutoff
radius $m$, with
\[
 10<a_1<b_1<m/10,\qquad 10m<a_2<b_2;
\]
\item closed Jordan disks $O\subset\operatorname{int}H$, with
$\overline B(0,b_2)\subset\operatorname{int}O$, and an
orientation-preserving band chart
\[
 \Phi:[-12,12]\times[-3,3]\longrightarrow\R^2.
\]
Writing its coordinates as $(u,v)$, membership in $O$ and $H$ within the
chart is exactly $|u|\le3$ and $|u|\le8$, respectively; membership in the
interiors is given by the strict inequalities.  The full section
$\Phi(\{0\}\times[-3,3])$ is strictly inside $B(0,a_1)$;
\item one positive integer rank for each of the two windows
\[
 I_B=[-.3,.3],\qquad I_W=[-1.8,-1.4].
\]
\end{enumerate}
No metric regularity of the Jordan disks or chart is imposed.
\end{definition}

Let $\nu$ be the ordinary full-plane colored Poisson law at spacing
$\eps$.  The state $\sigma$ of a placed plug is the colored point data in
$O$.  Set
\[
 \rho_{\eps,h}=h\sqrt{\eps/h}=\sqrt{h\eps}.
\]
An edge is \emph{retained} if it belongs to a Delaunay triangle with
circumradius at most $\rho_{\eps,h}$.  Whether such a triangle incident to
a vertex exists is determined by points within distance
$2\rho_{\eps,h}$ of that vertex.  Thus retained-edge queries in a compact
set use only a shrinking collar, on every configuration.  By
Lemma~\ref{lemma:mesh}, all edges and triangles of a fixed local system
are retained with probability tending to one.

For each color, take the connected components of the induced retained
graph whose vertex positions belong to
$\Phi([-2,2]\times I_{\rm color})$ and which meet both $u\le-1.5$ and
$u\ge1.5$.  Order them by the first site in chart lexicographic order.
The component of the specified rank, when present, is its \emph{root}.
These graph queries lie compactly inside $O$.  The test $L^+$ requires,
for each color, a retained path from its root to $u\ge6.7$, all of whose
vertices, including its ends, lie in $\Phi([1,7.5]\times I_{\rm color})$.
The test $L^-$ uses $u\le-6.7$ and
$\Phi([-7.5,-1]\times I_{\rm color})$.

\begin{definition}[Usable states and admission]\label{def:usable-state}
Fix positive thresholds $\beta,\beta'$.  The state event $G$ requires:
\begin{enumerate}[label=\textup{(\roman*)}]
\item simple surrounding retained site circuits $S_1,S_2$ in the two gate
annuli, with their edge segments in those annuli, enclosing the
corresponding inner balls and containing at most two black sites each;
\item both roots and the inequalities
$\nu(L^+\mid\sigma)\ge\beta$ and
$\nu(L^-\mid\sigma)\ge\beta$.
\end{enumerate}
The conditional probabilities use independent ordinary data on $O^c$;
in particular they average over the exterior point geometry as well as
its colors.  A type is \emph{admitted} when $\nu(G)\ge\beta'$, and its
state law is
\[
 \mu=\nu(\sigma\in\cdot\mid G).
\]
The finite normalized menu and thresholds will be fixed in
Section~\ref{sec:pins}.  Admission and $\mu$ may depend on $h,\eps$.
\end{definition}

Choose the two circuits by deterministic measurable rules on the state.
They are compactly contained in $O$.  The physical cutoff ball of a plug
centered at $z_s$ is the closed ball $B_s=\overline B(z_s,m_sh)$.

\subsection{Connectors and their conditional success probability}

\begin{definition}[Connectors and simple systems]\label{def:simple-system}
A directed connector $e:s\to t$ between same-color plugs with disjoint
obstacles is an oriented rectangular band, with transverse parameter
$v\in[-3,3]$, attached by its full parameterized ends to the positive
port $u_s=8$ of $H_s$ and the negative port $u_t=-8$ of $H_t$.  Apart from
these attachments it avoids both disks.  Only the charts with $|u|\le8$
are prescribed at its ends.

A \emph{simple system} is a cyclic list of at least two distinct plugs
and connectors.  Its closed obstacles are pairwise disjoint, its exterior
bands are pairwise disjoint and avoid all other obstacles, and the
sectors from section $0$ of one plug to section $0$ of the next form an
embedded annular band with the prescribed directed longitude and
transverse parameter.  Denote the whole closed sector of $e$ by
$\mathcal B_e$.
\end{definition}

We also allow an orientation-preserving ambient homeomorphism $g_s$ near
each plug, supported in a preset bounded local region, replacing $H_s$
and $\Phi_s$ by $g_sH_s$ and $g_s\Phi_s$.  The state region, roots and
tests remain unchanged.  Choose the allowed uniform displacement small
enough that $O_s\subset\operatorname{int}(g_sH_s)$, that section $0$ and
a full central slab still lie strictly inside the first gate, and that
positive and negative chart portions meeting either gate keep their
signs.  These conditions follow from the compact separation of portions
away from the central slab.  The tolerance is further reduced after any
local stub and layer margins used below have been fixed.

The success event $S_e$ asks for a retained root-to-root path of each
color in its own tube of $\mathcal B_e$, using the actual end charts:
\[
 -.65\le v\le .65\quad\text{for black},\qquad
 -2.2\le v\le-1\quad\text{for white}.
\]
Both tubes run longitudinally from $u_s=.5$ to $u_t=-.5$, and all path
vertices, including the ends, must lie there.  The exterior query sets
for different connectors are disjoint.  Indeed their whole sectors
outside the closed states are separated; at a shared plug the common
section is shielded inside $O$.  The tube footprints and joining seams
have room for the shrinking retained-geometry collars.

Figure~\ref{fig:plug-sector} shows the state, gates and cutoff within a
plug, together with the inset success tubes in its outgoing sector.

\begin{figure}[t]
\centering
\begin{tikzpicture}[x=1.12cm,y=.83cm,>=Stealth,font=\small]
 \fill[gray!5,rounded corners=9pt] (-2.35,-2.1) rectangle (2.35,2.1);
 \draw[thick,rounded corners=9pt] (-2.35,-2.1) rectangle (2.35,2.1);
 \fill[gray!12,rounded corners=6pt] (-1.32,-1.8) rectangle (1.32,1.8);
 \draw[dashed,rounded corners=6pt] (-1.32,-1.8) rectangle (1.32,1.8);
 \fill[black!16] (-2.65,-.195) rectangle (5.0,.195);
 \fill[blue!12] (-2.65,-.66) rectangle (5.0,-.30);
 \draw[gray] (-2.65,.9)--(5.0,.9);
 \draw[gray] (-2.65,-.9)--(5.0,-.9);
 \draw[thick] (0,0) ellipse [x radius=.45,y radius=1.16];
 \draw[densely dotted] (0,0) ellipse [x radius=.73,y radius=1.40];
 \draw[thick] (0,0) ellipse [x radius=1.04,y radius=1.65];
 \draw[dashed] (0,-.9)--(0,.9);
 \draw[very thick,->] (-2.65,.08)--(4.7,.08);
 \draw[blue!60!black,very thick,->] (-2.65,-.48)--(4.7,-.48);
 \foreach \xx in {-.449,.449,-1.039,1.039}
    \fill (\xx,.08) circle[radius=.042];
 \node[anchor=south west] at (1.50,2.12) {$H_s$};
 \node[anchor=south] at (0,2.60) {$O_s$: retained state};
 \draw[->] (0,2.55)--(0,1.80);
 \node[anchor=south east,fill=white,inner sep=1pt] at (-.18,1.12) {$S_1$};
 \node[anchor=south west,fill=white,inner sep=1pt] at (.69,1.43) {$S_2$};
 \node[anchor=east] at (-.80,-1.37) {$B_s$};
 \draw[->] (-.75,-1.34)--(-.44,-1.13);
 \node[anchor=north,fill=white,inner sep=1pt] at (0,-1.91) {section $0$};
 \draw[->] (0,-1.78)--(0,-.80);
 \node[anchor=south,fill=white,inner sep=1pt] at (2.35,1.0) {positive port $u_s=8$};
 \draw[->] (2.35,.96)--(2.35,.40);
 \node[anchor=south] at (3.77,1.48) {$\mathcal B_e$ toward $t$};
 \node[anchor=west] at (3.0,.50) {black tube};
 \node[anchor=west] at (3.0,-1.08) {white tube};
 \node[anchor=south] at (-1.90,2.18) {negative port};
 \draw[->] (-1.90,2.05)--(-2.35,.30);
\end{tikzpicture}
\caption{A topological schematic of a plug and its outgoing sector.
The two gate circuits have only two black contacts with the displayed
rim passage.  The cutoff lies between the gate annuli.  The state and
the whole sector are retained in comparisons; the success tubes are
strictly inset.  The shapes and radii in this schematic are not to scale.}
\label{fig:plug-sector}
\end{figure}

\begin{lemma}[Conditional bridge]\label{lemma:bridge}
For every fixed connector and admitted endpoint types, there is $c_e>0$
such that, for all sufficiently small $\eps$ and almost every pair of
usable endpoint states,
\begin{equation}\label{plug:success}
 \nu(S_e\mid\sigma_s,\sigma_t)\ge c_e.
\end{equation}
The constants are uniform for rigid copies and for buffered families
with a uniformly bounded number of fixed-shape crossing rectangles.
Moreover, on good continuation geometry, the white part of $S_e$ has a
positive quenched lower bound if the two white side tests each have
conditional probability at least $\beta$ given that geometry and the
inner states.
\end{lemma}

\begin{proof}
Write $A=L_s^+\cap L_t^-$ and $\sigma=(\sigma_s,\sigma_t)$.  The endpoint
exterior restrictions are independent, so $\nu(A\mid\sigma)\ge\beta^2$.
For each color, continue through its open layer by a finite connected
network of buffered crossing corridors.  At each end include a transverse
bar in the nominal slab $5<|u|<6$, crossing the entire corresponding
window with margin; join these bars through the layer.  All continuation
queries lie strictly outside the state disks.  A bar intersects every
successful side path, since that path traverses the slab through the
window.  Retention and buffers allow passage between cell and primal
crossings without changing this separation argument.

Here is the corridor construction in the form that will also be used
for local preparations.  Approximate a compact trace in an open corridor
by finitely many horizontal and vertical polygonal legs on a sufficiently
fine fixed square grid.  Put thin crossing rectangles along the legs,
with both horizontal and vertical crossings in successive overlap
squares, and add the same overlaps at branch points.  Crossing paths
therefore join into the prescribed network.  Slightly lengthening and
shrinking the rectangles leaves room to convert them into retained
primal paths.  Along a charted simple trace the steps may be chosen with
arbitrarily small longitudinal extent, so concatenation in order gives
any prescribed positive bound on longitudinal backtracking.

Let $E$ assert the retained-mesh and quenched rectangle lower bounds in
these finitely many continuation footprints.  Fixed collars keep its
point queries outside $O_s\cup O_t$, independently of the states.
Proposition~\ref{prop:probability-input} and Lemma~\ref{lemma:mesh} give
\[
 \nu(E^c\mid\sigma)\le\delta_\eps\longrightarrow0
\]
uniformly over the usable states.  No good-geometry requirement deep
inside an arbitrary fixed state is needed: its successful side paths
already use retained edges.

Condition now on all point geometry and on $\sigma$.  On sites queried
by black constraints order the bits toward black, and on the disjoint
set queried by white constraints order them toward white.  Both $A$ and
the continuation success $K$ are increasing in this one product order;
constraints of the same color may overlap.  The two sets of colors have
positive separation by the layer and chart-tolerance margins.  On $E$,
conditional positive association and the crossing input give
\[
 \nu(A\cap K\mid\eta,\sigma)
 \ge \nu(A\mid\eta,\sigma)\nu(K\mid\eta,\sigma)
 \ge q_e\nu(A\mid\eta,\sigma),
\]
where $q_e>0$ is fixed.  Since $A\cap K\cap E\subset S_e$, integration
before subtraction yields
\[
 \nu(S_e\mid\sigma)
 \ge q_e\nu(A\cap E\mid\sigma)
 \ge q_e(\beta^2-\delta_\eps).
\]
This proves \eqref{plug:success}.  It does not infer any quenched black
side propensity from an annealed one.  If the two white side propensities
are themselves bounded below after fixing geometry, the white-only
version of the same product argument gives $q_e^W\beta^2$ without this
annealed integration.  For fixed normalized templates, finitely many
buffered rectangle chains give common positive constants by scaling;
rigid motions preserve the laws.
\end{proof}

\subsection{Cycle experiments and exact path agreement}

\begin{definition}[Cycle experiment]\label{def:cycle-experiment}
For a simple system, sample the states independently with their laws
$\mu_s$.  Given the states, sample the exterior regions of the connectors
independently, each conditioned on its success $S_e$, and sample the
remaining plane data ordinarily.  The exterior region of $e$ may be the
whole sector $\mathcal B_e$ with the closed endpoint state disks removed.
\end{definition}

All queries of $S_e$ are then in its region or its endpoint states, for
small $\eps$.  The cycle law has density
\begin{equation}\label{plug:cycle-density}
 \prod_s\frac{\ind_{G_s}}{\nu(G_s)}
 \prod_{e:s\to t}
 \frac{\ind_{S_e}}{\nu(S_e\mid\sigma_s,\sigma_t)}
\end{equation}
relative to $\nu$ on these data; unused data remain free.  Thus it is
boundedly dominated by $\nu$ for a fixed system.  In particular, the
marginal of two endpoint states and one sector is exactly their ordinary
law restricted by the three factors for $G_s,G_t,S_e$.  Deterministic
Borel restrictions suffice throughout: a Jordan boundary need not be
smooth or have zero area.

\begin{lemma}[A confined rim and its gate passages]\label{lemma:cycle-rim}
On a cycle experiment, outside an event of probability tending to zero,
there is a rim $P$ of the path color, oriented in the positive longitude,
contained in the layer $-2.35<v<.8$.  At every selected gate circuit it
has exactly two contacts, crossing from the negative side into the
interior and then out through the positive side.  Its outside path from
the first-gate exit at $s$ to the first-gate entry at $t$ lies in
$\mathcal B_e$.  The inner gate paths occur in the prescribed cyclic
order.
\end{lemma}

\begin{proof}
The root components join negative to positive side vertices within their
central windows.  Concatenate these paths with the successful sector
paths.  For both colors one obtains a monochromatic closed walk of
longitudinal index one, respectively in the disjoint layers
$-.75<v<.75$ and $-2.3<v<-.9$.  This index follows by lifting the
longitude: root passages cross the central slabs, and sector passages go
from one positive slab to the next negative one without passing the end
section cuts.  Arbitrary local backtracking is harmless.  Fine mesh
keeps the edge segments in the stated layers.  Each walk contains an
essential simple circuit.  Lemma~\ref{lemma:sandwich} supplies a black
rim between them.  The farther subbands on either side lie in the two
exterior sides of this sandwich, so the chosen rim is confined to
$-2.35<v<.8$.

This rim visits both the inside and outside of either gate: it crosses
section $0$ deep inside and also reaches the other plugs.  A black rim
can meet a gate with at most two black sites only at those sites.  By
simplicity it must meet exactly two, once each, and cross at both.  A
tangency would not change side, leaving too few contacts.  Following a
gate edge between the two contacts is also impossible: the complementary
rim arc would then have to visit both sides without another contact.
Thus the incoming and outgoing edges are strictly in opposite gate fans.

All visits to the central section belong to the single inner gate arc.
That arc lies in the end chart, since the gate and its bounded inside lie
in $O$ and all other bands avoid its obstacle.  The full central slab is
inside the first gate, and compact positive and negative chart portions
meeting the gates are separated.  The total positive longitudinal index
one therefore forces negative entry and positive exit.  This uses
ordered passage, or equivalently lifted longitude, without requiring a
transversal intersection with section $0$.  Outside its inner first-gate
arc the rim cannot cross that section.  Consequently the outside arc
from $s$ to the next $t$ stays in their intervening sector.  Since each
inner arc contains all visits to its own cut and these chart slabs are
disjoint, the inner arcs have the asserted cyclic order.
\end{proof}

The following formulation isolates the precise deterministic comparison
that will also apply to rims in different configurations.

\begin{lemma}[Two-graph gate consistency]\label{lemma:gate-consistency}
Let $P,P'$ be oriented simple monochromatic rims in their respective
embedded triangulations.  Suppose their directed subpaths $\gamma,
\gamma'$ have common distinct initial and terminal vertices $v,w$ and
satisfy these conditions:
\begin{enumerate}[label=\textup{(\roman*)}]
\item each \emph{complete compared subpath} is present and monochromatic
in the other graph, with the same embedding;
\item at $v$ there is a shared gate circuit, both incoming rim rays lie
strictly in one gate fan, and both outgoing rays lie strictly in the
opposite fan;
\item the rim orientations and direction of crossing that gate agree.
\end{enumerate}
Then $\gamma=\gamma'$ as directed edge paths.  Each rim needs its
no-right-bypass property only in its own graph.
\end{lemma}

\begin{proof}
Order rays in the outgoing fan counterclockwise between its two boundary
rays.  For any incoming ray in the opposite open fan, the
counterclockwise sector from that ray to an outgoing ray contains the
same initial portion of the outgoing fan.  Hence, if the two outgoing
rays differ, one enters the other's open right sector, independently of
the differing incoming rays.  Suppose, for instance, that $\gamma'$
leaves to the right of $P$.

Follow $\gamma'$ to its first subsequent contact with the \emph{whole}
rim $P$, not merely with $\gamma$.  Such a contact exists no later than
$w$.  It is not $v$, because $\gamma'$ is simple.  As $\gamma'$ is a path
in $P$'s graph, planarity makes the first contact a vertex contact; a
shared edge is first reached at its endpoint.  Before this contact the
path is connected and disjoint from $P$, so it stays in the open right
component into which it departed.  It is therefore a monochromatic
right bypass between distinct vertices of $P$, contradicting
Lemma~\ref{lemma:no-bypass}.  An earlier return on the complementary
portion of $P$ gives exactly the same contradiction.  The other case
uses $P'$ in its own graph.

Thus the first edges agree.  At any later first divergence the incoming
edge is common, and cyclic order from that incoming ray again puts one
outgoing edge into the other's open right sector.  The same whole-rim
first-return argument rules this out.  The paths agree to $w$.
\end{proof}

\begin{lemma}[Deterministic collars for the compared paths]
\label{lemma:fixed-collar}
For a fixed cycle pattern, every possible confined experiment rim has
its inner second-gate path in a fixed compact subset of its state disk.
Every possible first-exit-to-next-first-entry path lies in a fixed
compact subset of
\[
 D_e=O_s\cup O_t\cup\mathcal B_e.
\]
Consequently, if two configurations retain these data, their complete
incident triangulations along the compared paths agree on the mesh-good
event, for all sufficiently small $\eps$.
\end{lemma}

\begin{proof}
The second gate is in its prescribed outer ball, compactly inside $O$;
its bounded inside is in the same ball.  This gives one deterministic
envelope, also containing the gate neighborhoods.  For an outside path
use the closed inset footprint $-2.35\le v\le.8$ in the whole sector.
It has positive lateral margins.  The chart--connector seams are
interior product-band seams.  At the two longitudinal ends the full
section and a central slab are compactly inside the endpoint states,
which shield the cuts.  Hence this footprint is compactly contained in
$\operatorname{int}D_e$.  Its collar is fixed before selecting a rim;
there is no requirement that the rim follow a particular witness of
$S_e$.

If a triangle is incident to a vertex in such an envelope and its
circumradius is at most $\rho_{\eps,h}$, its whole circumdisk is within
$2\rho_{\eps,h}$ of that envelope.  For small $\eps$ the disk lies in
the retained region, so its vertices and empty-circumdisk condition are
unchanged.  Mesh control in \emph{both} configurations excludes additional
large incident triangles on either side.  Thus all incident triangles,
not only the edges of a preselected path, agree.  The required mesh
probabilities follow from ordinary control and the fixed density bounds
for the experiments.
\end{proof}

For coupled cycle experiments sharing a state, choose the same
deterministic gate circuits from that state.  By
Lemma~\ref{lemma:cycle-rim}, both rims visit the inside and outside of
each chosen gate and can meet it only at its at most two path-color
vertices.  They therefore use the same two distinct vertices; negative
entry and positive exit give the same ordering of these contacts.  Thus
the inner second-gate paths have common ordered endpoints, and the
outside paths from the first-gate exit at $s$ to the first-gate entry at
$t$ do also when both endpoint states are shared; these endpoints are
distinct because the state disks are disjoint.  At the start of an
inner path, both incoming rays are in the exterior gate fan and both
outgoing rays in the interior fan; at the start of an outside path these
roles are reversed.  All these inclusions are strict.  Sharing the state
preserves the colored collar of each complete inner second-gate path;
sharing both endpoint states and the whole sector does the same for
each complete outside path.  On the mesh-good event,
Lemma~\ref{lemma:fixed-collar} therefore puts each compared path in the
other graph with the same colors and preserves every incident triangle.
Lemma~\ref{lemma:gate-consistency} now identifies the directed paths.
Their right incident triangles and local tagged sides consequently
agree as well.  The allowed chart changes preserve the sign and collar
margins used in this argument.

\subsection{Removing the exterior randomness}

On a cycle experiment choose any measurable confined rim as above.
For $X=Q,V$ define
\begin{equation}\label{plug:masses}
 \widehat m_s=X_P(\ind_{B_s}),\qquad
 \widehat j_e=X_P\bigl(\ind_{\mathcal B_e\setminus(B_s\cup B_t)}\bigr),
 \qquad e:s\to t.
\end{equation}
These are barycenter masks.  Fix disjoint Borel conventions at sector
boundaries; the rim is buffered away from lateral boundaries and the
section cuts are inside cutoff balls.  One may set the masses to zero
on the exceptional geometry event.  The deterministic-mask energy
bound \eqref{fields:energy}, ordinary triangle moments and bounded cycle
density give finite $L^2$ bounds for each fixed pattern.  For bounded
normalized short templates with common density bounds these norms are
at most $C h$.

Every cutoff tag comes from the inner path through the \emph{second}
gates, because the cutoff is buffered inside that gate.  The sector
mass is precisely the contribution outside both cutoff balls along the
path from the \emph{first} exit at $s$ to the first entry at $t$:
the inner first-gate paths and cuts lie safely inside the balls.  Thus,
on the mesh-good event,
\[
 \sum_{e:s\to t}(\widehat m_s+\widehat j_e)=X_P(1).
\]

\newpage
\begin{proposition}[The state cochain]\label{prop:state-cochain}
For every fixed connector, including its actual end charts, there are
functions $i_e^\eps(\sigma_s,\sigma_t)$ such that
\begin{equation}\label{plug:cycle}
 \begin{aligned}
 \widehat m_s+\widehat j_e-i_e^\eps&=o_{L^1}(1)
       &&\text{in every cycle experiment containing }e,\\
 \sum_e i_e^\eps&=o_{L^1}(1)
       &&\text{on the independent state product of every simple system}.
 \end{aligned}
\end{equation}
The same functions work in all completions.  For a fixed collection of
normalized short templates admitting uniform short completions,
\[
 \|i_e^\eps\|_{L^2(\mu_s\otimes\mu_t)}\le C h
\]
for all sufficiently small $\eps$, after increasing a common constant.
\end{proposition}

\begin{proof}
For each type, fix a short unperturbed reference completion, for example
using a sufficiently separated copy.  Define $m_s^\eps(\sigma_s)$ as
the conditional mean of its cutoff mass given $\sigma_s$.  Couple any
other experiment to this completion sharing only that state and sampling
the remaining data conditionally independently.  Inner second-gate
consistency identifies the cutoff mass on their common mesh-good event.
Their $L^2$ bounds show that their difference is $o_{L^1}(1)$; indeed
$\EE[|Y|\ind_E]\le\|Y\|_2\PP(E)^{1/2}$ on a bad event $E$.
Using two copies also in the reference completion and conditional Jensen
shows that its mass is $o_{L^1}(1)$ from its own conditional mean.
Consequently $\widehat m_s=m_s^\eps+o_{L^1}(1)$ in every completion.

Next fix one two-plug completion of $e$ and condition its sector mass on
$\sigma_s,\sigma_t$ and the exterior sector datum $\xi_e$.  Call the
result $j_e^\eps(\sigma_s,\sigma_t,\xi_e)$.  These retained data have
the same marginal in every completion, by \eqref{plug:cycle-density}.
Coupling and outside first-gate consistency, again with a second copy
for the defining completion, give
$\widehat j_e=j_e^\eps+o_{L^1}(1)$ in every experiment containing $e$.

In a fixed cycle, conditional on its complete state vector, the data
$\xi_e$ are independent, each with the conditioned law determined by its
two endpoint states.  The closed estimate \eqref{fields:closed} and the
cycle density show that
\[
 \sum_e(m_s^\eps+j_e^\eps)=o_{L^1}(1).
\]
For $Q$ the selected rim is contained in a fixed box and has a fixed
positive minimum diameter, since it visits distinct plugs; for $V$ its
complete tangent sum is zero.  Now resample one $\xi_e$ independently
given the states, leaving all other sectors fixed.  Subtracting the two
small cycle sums gives $j_e^\eps-(j_e^\eps)'=o_{L^1}(1)$.  Conditional
Jensen therefore gives
\[
 j_e^\eps-
 \EE[j_e^\eps\mid\sigma_s,\sigma_t]=o_{L^1}(1).
\]
The conditional law used in this last expectation depends only on those
two states.  Set
\[
 i_e^\eps=m_s^\eps+
         \EE[j_e^\eps\mid\sigma_s,\sigma_t].
\]
This proves both statements in \eqref{plug:cycle}.  The defining short
completions and Jensen give the stated $L^2$ bounds.  All replacements
use only bounded second moments and high-probability path equality;
uniform integrability of their squares is unnecessary.  The choices
are measurable because the locally finite graphs have countably many
finite paths and walls, and different choices agree through the same
path comparison.
\end{proof}

\subsection{Common closures, collars and changes of charts}

\begin{lemma}[Comparison by a common closure]\label{lemma:common-closure}
Suppose two directed paths of connectors from $s$ to $t$ admit a common
complementary connector $c:t\to s$, making each a simple system.  Their
sums of cochain values differ by $o_{L^1}(1)$ on the product of all states,
with shared states identified.  Either path may be a single connector.
For fixed end disks and charts, a connector value therefore depends,
up to $o_{L^1}(1)$, only on the endpoint-fixed deformation class of its
exterior simple centerline described below.
\end{lemma}

\begin{proof}
Apply \eqref{plug:cycle} to the two cycles and subtract the common
$i_c^\eps$.  Adding independent unused states does not change an $L^1$
error.  A complementary connector always exists: two disjoint closed
Jordan disks joined by the attached rectangular band form a Jordan disk,
with the unused ports on its boundary.  In its exterior, the
Jordan--Schoenflies theorem gives another handle band joining those
ports.  The transverse parametrizations match because departure and
arrival sections have opposite induced boundary orders.  The remaining
deformation assertion follows from the collar construction and finite
chaining below.
\end{proof}

\begin{lemma}[A band with prescribed full attachments]\label{lemma:band-collar}
Let a simple exterior arc join the middle points of two prescribed port
sections.  Suppose it continues for a short distance along the given
product collars at its ends and otherwise avoids the closed end disks.
Then a rectangular band about it, with those full parameterized
attachments, fits in any neighborhood of the arc and attachments within
the exterior.  It can in addition avoid any fixed disjoint compact
routes or unused obstacles.
\end{lemma}

\begin{proof}
Extend a short piece of each given end collar and taper its transverse
width to be very small at the far tip.  The full port is usable initially:
the compact middle arc is separated from it, while the given product
collar handles its immediate neighborhood.  Straighten a slightly
extended middle arc by an orientation-preserving plane arc chart, with
its left side above a horizontal segment.

From the upper and lower corners of each narrow collar tip, follow
constant small signed transverse coordinates in that end collar for a
short additional distance.  In the straightened coordinates these side
paths lie on the corresponding sides of the midline and progress along
it with arbitrarily small error.  Stop at their first hits of prescribed
vertical lines, chosen in order at the two ends.  Join the upper stopped
paths, and separately the lower ones, by height graphs strictly on their
side of the middle segment and as close to it as desired.  The compact
middle interval is separated from the initial collar pieces and disks,
so these joins fit in the specified neighborhood and avoid all the
specified compact sets.  First-hit stopping keeps the side paths simple.

Each joined side, the appropriate half-tip segments and the centerline
bound a Jordan rectangle.  The two rectangles lie on opposite sides in
the straightened coordinates and may be taken thin there.  The initial
disk-with-collar interiors abut their tips on the other side, so they
cannot cross either boundary curve.  Fill the rectangles by
Jordan--Schoenflies, prescribing the centerline and boundary
parametrizations, and glue them to the end collars.  This gives the
required band with the full original port attachments.
\end{proof}

For completeness, the deformation in
Lemma~\ref{lemma:common-closure} is implemented as follows.  Normalize
the disjoint Jordan disks to round holes by an ambient homeomorphism.
Inside an existing exterior band replace its centerline by a smooth
proper arc, radial near its endpoints, retaining the same complementary
connector.  Indeed an attaching section has a relative exterior
neighborhood; polygonal approximation in this open set, loop erasure
and smoothing give a simple arc with the prescribed short radial germs.
The full port attachments are recovered by
Lemma~\ref{lemma:band-collar}.

Consider now an endpoint-fixed deformation through smooth proper arcs,
transverse at the hole boundaries.  Product end collars are available
throughout, adjusting their angular coordinate with the radial level
near each endpoint.  At a parameter value choose a narrow connector and
a complementary band disjoint from it.  By compact separation that
complement also avoids all nearby centerlines; the collar lemma supplies
connectors for them with the same full attachments.  At overlapping
parameter neighborhoods, narrow a connector to avoid both choices of
complement.  A finite subcover of the parameter interval and the
common-closure comparison give the $o_{L^1}(1)$ equality along the whole
deformation.  Constants may depend on this finite chain, which is fixed
before $\eps\downarrow0$.

\begin{lemma}[Stability on buffered local templates]\label{lemma:chart-stability}
Fix a buffered two-plug system $e_0,c_0$.  Suppose its parts in the preset
neighborhoods where charts may move are fixed local stubs, and no other
part of the system returns to these neighborhoods.  For sufficiently
small supported changes $g$ of the two end charts and obstacles, let
$e_g,c_g$ be the corresponding images.  Then
\[
 i_{e_g}^\eps=i_{e_0}^\eps+o_{L^1}(1)
       \quad\text{on }\mu_s\otimes\mu_t.
\]
The required tolerance depends only on the compact local stub and layer
margins in the moved regions.
\end{lemma}

\begin{proof}
Consider the mixed experiment with successes $S_{e_g}$ and $S_{c_0}$,
with the nominal states and roots.  Their exterior queries are still
separated, and Lemma~\ref{lemma:bridge} applies.  All allowed tubes lie
in slightly widened, separated layers of the nominal annular band,
with the same longitudinal order.  This follows by the local tolerance
where anything moved; elsewhere the system is unchanged.  Thus the mixed
experiment has the same confined-rim and gate conclusions, and we may
use deterministic nominal masks for its masses.

Compare its $e$ contribution with the pure $e_g,c_g$ experiment by
sharing the states and the exterior part of the \emph{nominal} $e_0$
sector.  Every possible pulled-forward inset $e_g$ path outside the
states, together with all queries of $S_{e_g}$, is buffered in that
shared region.  The queries of both $S_{c_g}$ and $S_{c_0}$ avoid it.
Consequently the shared marginal in both experiments is the ordinary
one conditioned on $S_{e_g}$ at the two sampled states: the other
success factor integrates to one given the states and affects only
other exterior data.  This also follows directly from the product
density \eqref{plug:cycle-density}, irrespective of how unused points
are assigned to the sampling regions.

The common collar and the two gate comparisons identify the first-gate
path outside the cutoff balls and the starting cutoff contribution.
Thus the extracted $e$ value in the mixed experiment is $i_{e_g}^\eps$
up to $o_{L^1}(1)$.  The same argument identifies its other contribution
with $i_{c_0}^\eps$.  Closing the mixed rim therefore gives
$i_{e_g}^\eps+i_{c_0}^\eps=o_{L^1}(1)$, whereas the nominal system gives
$i_{e_0}^\eps+i_{c_0}^\eps=o_{L^1}(1)$.  Subtraction proves the claim.
\end{proof}

The restriction to buffered local templates is deliberate.  An
arbitrarily thin fitted connector with changed charts is first deformed
with those changed charts held fixed.  Only its standard buffered
representative is then compared to the nominal charts by
Lemma~\ref{lemma:chart-stability}.  A narrow fitted band's conditioning
constant is never used as the uniform constant of a short-template
moment estimate.
